\lesson{1}{Nov 01 2021 Mon (18:33:36)}{Valence Electrons}{Unit 3}

\subsubsection*{Factors that Influence Chemical Bonding}

Compounds are held together by \bf{Electrostatic Forces} between atoms in a \bf{Chemical Bond}. All chemical bonds involve \bf{Valence Electrons}, but bonds are classified by the way electrons are distributed within them. Here are the two types of bonds:

\paragraph{Ionic Bonds} An \bf{Ionic Bond} is a chemical bond that results from the \bf{Electrostatic Attraction} between positive ($+$) and negative ($-$) ions. This means that in an ionic bond, electrons are given up by one atom and gained by another atom, and then those atoms are attracted to each other.

\paragraph{Covalent Bonds} In a \bf{Covalent Bond}, electrons are shared between two atoms, neither atom completely gaining or losing electrons. Because atoms in \bf{Covalent Bonds} overlap their outer energy levels, both atoms \it{"Own"} the shared electrons.

In general, atoms of \bf{nonmetals} form \bf{covalent bonds} with each other, and \bf{metal atoms} form \bf{ionic bonds} with \bf{nonmetal atoms}. There are many exceptions to this generalization, but it still helps us predict which elements will bond together and what type of bond they will form.

\subsubsection*{Valence Electrons Identified in Electron Configurations}

The most important electrons in an atom are the ones in its \bf{Outermost} or \bf{valence shell}. The \bf{outermost shell} is the highest energy level that contains electrons in an atom. The valence electrons within the \bf{outermost shell} can be removed from one atom and given to another atom. These electrons can also be shared with atoms of other elements. This is how atoms undergo \bf{chemical reactions} or bond with other atoms. For the element krypton, the highest energy level used is level $4$. There are two occupied \bf{subshells} in level $4$—\bf{subshells} $4s$ and $4p$. The $s$ \bf{subshell} has $2$ electrons, and the $p$ \bf{subshell} has $6$ electrons. This totals to $8$ \bf{valence electrons} for krypton.

\subsubsection*{Using the Periodic Table to determine an Element's Valence Electrons}

The \bf{Octet Rule} states that atoms for elements in the main element groups prefer to have up to $8$ \bf{valence electrons}. This fits how the periodic table is organized, with $8$ tall columns and several shorter columns for \bf{transition metals}. If the \bf{valence electrons} equal $8$, the element is inert, which means it doesn't react with other substances. If an atom of an element has fewer than $8$ electrons, it may seek to form chemical reactions with other atoms to get a full outer shell of valence electrons.

\paragraph{Columns $1$ and $2$} All the elements in column $1$ have $1$ valence electron in their \bf{outermost shell}. For column $2$, it's the same idea. Atoms of elements in this column have $2$ \bf{valence electrons}.

\paragraph{Transition Metals} The transition metals follow unique rules, so their columns cannot directly tell us about their total \bf{valence electrons}. Scientists determine their valence electrons in other ways.

\paragraph{Columns $3$ to $7$} Just like columns one and two, taller columns three through seven indicate the number of valence electrons for the elements within them. For example, elements in column three have three electrons.

\paragraph{Column $8$} With the exception of helium, which only needs $2$ electrons to fill its \bf{outermost shell}, all other elements in column $8$ have $8$ \bf{valence electrons}.

\subsubsection*{Element's Valence Electrons relate to its Chemical Reactivity}

Valence electrons are often found in the s and p sublevels of the outermost energy level in an atom. Each subshell has its own unique shape. Subshells d and f have some challenges to bonding that the s and p subshells do not have. That is why d and f electrons are not usually included in the valence and are rarely involved in chemical reactions.

Empty spaces in the subshells of the outermost shell also affect reactivity. Atoms prefer eight valence electrons. Therefore, they will bond with other substances to gain a full eight. Some atoms may gain electrons to make eight, while others can give electrons away. Let's see how that works.

\subsubsection*{Valence Electrons affect the Bonding Properties of Elements}

Elements with the same number of valence electrons show similar chemical behavior, which is why they are grouped together on the periodic table. These elements also follow similar periodic trends when it comes to ionization energy and electronegativity.

Chemical bonds can range in strength and properties between completely ionic and completely covalent, depending on how strongly the atoms attract the electrons. Comparing their electronegativity values helps chemists estimate the degree of ionic or covalent properties of various chemical bonds. The greater the difference in their electronegativity values, the more ionic the chemical bond between two atoms.

\subsubsection*{Shorthand Notation to Illustrate Valence Electrons}

Full \bf{Electron Configurations} can get very long when atoms have large atomic numbers. Sometimes, the chemist is only concerned with an atom's valence electrons, since those are the electrons that are more involved with bonding and other chemical properties. For this cases, scientists have created a special notation, called \bf{Gas Notation}, for abbreviating the electron configuration.

For example, for \bf{Chlorine (Cl)}, the electron configuration is $1s^22s^22p^63s^23p^5$. Using the \bf{Gas Notation}, we can write the electron configuration as \bf{[Ne] $3s^23p^5$}

\begin{center}
    How to Determine the Electron Configuration for Mercury (Hg)
    
    1. Write the symbol of the last element from the previous row in brackets. [Xe]
    2. Mercury's electron configuration is $1s^22s^22p^63s^23p^64s^23d^{10}4p^65s^25p^66s^24f^{14}5d^{14}$. Write the rest of the element's notation to complete that element's shorthand notation: \\
    
    \begin{align}
        \rm{[Xe]} 6s^24f^{14}5d^{10}
    \end{align}
\end{center}

\newpage
